\documentclass{article}
\usepackage[utf8]{inputenc}
\usepackage[T1]{fontenc}
\usepackage{graphicx}
\usepackage{algorithm}
\usepackage{algpseudocode}

\usepackage{listings}
\usepackage{xcolor}
\usepackage[margin=2.5cm]{geometry}
\usepackage{float}

\usepackage{amssymb} %Ajoute des symboles mathématiques comme R, Z, forall, exists, etc....
\usepackage{amsmath}
\usepackage{amsthm}


\lstdefinestyle{monstyle}{
    backgroundcolor=\color{white},
    commentstyle=\color{gray},
    keywordstyle=\color{blue},
    numberstyle=\tiny\color{gray},
    basicstyle=\ttfamily\small,
    breaklines=true,
    numbers=left,
    numbersep=5pt,
    frame=single,
}
\lstset{style=monstyle}

\title{\textbf{Rapport de TD n°3 : Initiation à l'algorithmique}}
\author{Derrez Lorenzo}

\begin{document}
\maketitle

\section*{\textbf{Exercice 0 :}}

\section*{\textbf{Exercice 1 :}}

\noindent \textbf{Hypothèses :} 1 opération = $10^{-9}$ s (une machine à 1 milliard d'opérations par seconde). Pour une entrée de $n = 96$ octets, un algorithme linéaire effectue $n$ opérations et un algorithme exponentiel effectue $2^n$ opérations.

\medskip
\noindent Pour le $1^{er}$ cas :
\begin{itemize}
    \item Complexité linéaire $\implies$ proportionnel à la taille en octets de l'entrée.
    \item Après calcul on obtient en temps linéaire $96 \times 10^{-9}\ \text{s} \approx 0{,}1\ \mu\text{s}$.
\end{itemize}

\medskip
\noindent Pour le $2^{e}$ cas :
\begin{itemize}
    \item Complexité exponentielle $\implies$ $2^{96}$ opérations.
    \item Calcul mental avec $2^{10} \approx 10^3$ : $2^{96} = 2^6 \times (2^{10})^9 \approx 64 \times 10^{27} \approx 6 \times 10^{28}$ opérations.
    \item Temps : $2^{96} \times 10^{-9}\ \text{s} \approx 8 \times 10^{19}\ \text{s}$.
    \item Avec $1\ \text{an} \approx 3 \times 10^{7}\ \text{s}$, on obtient en temps exponentiel environ $2{,}5 \times 10^{12}$ ans, soit environ 2\,500 milliards d'années.
\end{itemize}

\section*{\textbf{Exercice 2 :}}
\subsection*{Question 1}
\begin{lstlisting}[language=Python, caption={Complétion de la fonction exo 2}]
def produit (A : Matrix[integer], B : Matrix[integer]) -> Matrix[integer] :
    (h,w) = A;shape
    C = newMatrix(h, h, 0)
    for i in range(h) :
        for j in range(h) :
            for k in range(h) :
                C[i, j] = C[i, j] + A[i, k] * B[k, j]
    return C 
\end{lstlisting}

\subsection*{Question 2}
\text{Sur les entrées :}
\qquad A, B -> $h^2$ (car w = h par hypothèse)
\\
\qquad On obtient : $O(h^2)$

\subsection*{Question 3}
\qquad 3 boucles imbriquées $\implies$ $O(h^3)$

\subsection*{Question 4}
\qquad  C'est cubique en h. 

\subsection*{Question 5}
\qquad C occupe une place $h^2$ $\implies O(h^2)$ (ici i, j et k sont en $O(1)$)

\subsection*{Question 6}
\qquad C'est quadratique en h. 

\subsection*{Question 7}
\qquad On pourrait réduire la complexité en espace en modifiant directement la valeur de la matrice A ligne par ligne plutôt qu'en créant une nouvelle matrice C.

\subsection*{Question 8}
\qquad En cherchant sur WWW, il est dit qu'il existe un autre algorithme, celui de \textbf{Strassen}. Ce dernier implique : $O(h^{2,807})$.


\section*{\textbf{Exercice 3 :}}
\subsection*{Question 1 : Arbre des appels}


\subsection*{Question 2}
-- Nombre d'appels récursifs : 
\\
-- Puis, on obtient le tableau suivant selon les types : 

\section*{\textbf{Exercice 4 :}}

\section*{\textbf{Exercice 5 :}}

\subsection*{Question 1}
\begin{lstlisting}[language=Python, caption={fonction \textit{longueurMaximalePlateau}}]
def longueurMaximalePlateau(T: Array[integer]) -> integer :
    int n = T.size
    int max = 0
    for i in range(0, n) :
        int j = i + 1
        while (j <= n and estPlateau(i, j, T)) :
            j = j + 1
        # T[i:j-1] est le plus grand plateau commencant en i
        if (j - 1 - i > max) :
            max = j - 1 - i
    return max
\end{lstlisting}

\begin{itemize}
    \item \textbf{Complexité en temps :} 
    \subitem 3 imbrications en lien avec n la taille du tableau $\implies$ $O(n^3)$ (car \textit{estPlateau} est linéaire en n)
    \item \textbf{Complexité en espace :}
    \subitem toutes les opérations sont constantes $\implies O(1)$.
\end{itemize}

\subsection*{Question 2}
\begin{itemize}
    \item Il faut un exemple avec une chaine maximale qui n'est pas au début, et qui n'est pas à la fin non plus.
    \item On peut prendre comme exemple : $T = [1, 1, 3, 3, 3, 2, 2, 7]$
    \item Un taille de tableau égale à 5 serait surement possible, par exemple : $T = [1, 3, 3, 3, 7]$
\end{itemize}

\subsection*{Question 3}
\begin{lstlisting}[language=Python, caption={Complétion du modèle de la fonction \textit{IMP}}]
    def lmp (T : Array[integer]) -> integer : 
        int lm = 0
        int current = 0
        for i in range(0, T.size) : 
            if i > 0 and T[i] == T[i-1] : 
                current = current + 1 
            else : 
                current = 1
            if current > lm : 
                lm = current
        return lm 
\end{lstlisting}

\end{document}